% Chapter introduction: why the topic arose, who first asked, what it answers.
% STATUS: rewritten 2026-11-05
% Sources read in full: Filippov 2000 (pp. 354--358), McCrimmon 2004 (pp. 52,
% 160, 185), Davenport 1972 (p. 565), Loday and Pirashvili 1993 (pp. 139--141),
% Ayupov, Kudaybergenov, Omirov and Zhao 2019 (pp. 3, 11), Biyogmam and
% Tcheka 2021 (p. 2193), Musson 2012 (pp. 2--3), Kac 1977 (p. 1375),
% Hartwig, Larsson and Silvestrov 2006 (p. 18), Jin and Li 2008 (pp. 1399--1400),
% Li 2018 (pp. 657--659), Ferreira, Kaygorodov and Lopatkin 2021 (pp. 1, 4--5),
% Fan, Hong and Su 2019 (pp. 2--3), Jiménez-Gestal and Pérez-Izquierdo 2003,
% Kaygorodov 2014 (Theorem 4), Baez 2002 (pp. 148, 159).
The chapters so far stayed inside associative rings, Lie algebras, and
operator algebras. The derivation algebra of a Lie algebra is again a Lie
algebra, and many classical results about derivations use associativity of
the underlying product. In this chapter we ask what survives when the
product is no longer associative nor antisymmetric. The classes include
the Jordan product of hermitian matrices, the octonionic product, the
Leibniz identity that Loday isolated in 1993, and the brackets of
superalgebras. They also include the twisted Jacobi rules of Hom-Lie
algebras and Lie conformal algebras.

Some of these questions are older than the pages below may suggest.
Jacobson realised in the 1950s that the exceptional Lie algebras of types
$G_2$ and $F_4$ arise as derivation algebras of the octonions and of the
Albert algebra \cite[p.~52]{McCrimmon2004Taste}. Loday introduced Leibniz
algebras in 1993 to explain why a certain chain complex in cyclic homology
is closed \cite[p.~139]{LodayPirashvili1993}. Hartwig, Larsson and
Silvestrov isolated a twisted Jacobi identity in 2003 to describe
deformations of the Witt algebra
\cite[p.~1]{HartwigLarssonSilvestrov2006}. Each line has its own reason
for existing and its own derivative-like linear maps.

What we collect below is the literature on derivations in these wider
classes, with one result from each. We ask in every section: what is a
derivation here, does the analogue of ``every derivation of a simple
finite-dimensional algebra is inner'' hold, and what does the derivation
algebra compute?

Section~\ref{sec:nonassociative:jordan-algebras} works with Jordan
algebras proper. The inner derivations of a Jordan algebra are the maps
$D_{x,y}=4[L_x,L_y]=[V_x,V_y]$. They span a Lie ideal of $\mathrm{Der}(J)$
\cite[Problem~3.6, p.~185]{McCrimmon2004Taste}. The $27$-dimensional
Albert algebra of hermitian octonionic $3\times 3$ matrices has derivation
algebra $F_4$ \cite[p.~52]{McCrimmon2004Taste}.
Section~\ref{sec:nonassociative:alternative} records the definition of Artin
of an alternative algebra and the theorem of Davenport that derivations of a
semisimple Malcev algebra in characteristic $\neq 2,3$ are inner
\cite[p.~148]{Baez2002Octonions}, \cite[p.~565]{Davenport1972Malcev}.
Section~\ref{sec:nonassociative:leibniz} opens the definition of Loday and
gives the derivation structure of complex semisimple Leibniz algebras
\cite[Definition~(1.1), p.~141]{LodayPirashvili1993},
\cite[Theorem~4.5]{AyupovKudaybergenovOmirovZhao2019},
\cite[Proposition~3.5, p.~2193]{BiyogmamTcheka2021Outer}. At dimension
three the description of the derivation algebras is complete
\cite[Theorem~1.1, p.~2]{KurdachenkoPypkaSemko2026}.
Section~\ref{sec:nonassociative:superalgebras} states the definition of Musson
of a superderivation and ties it to the classification of Kac of simple
Lie superalgebras \cite[Lemma~1.1.2, p.~2]{Musson2012},
\cite[p.~1375]{Kac1977Classification}. It records that a two-step nilpotent
Lie superalgebra carries local superderivations that are not
superderivations \cite[Theorems~3.2 and 3.5, pp.~7, 12]{ChiZhangMengTang2026}.
Section~\ref{sec:nonassociative:delta-quasi-anti} treats the
$\delta$-derivations of alternative and Malcev algebras (Filippov 2000)
and of simple Jordan superalgebras (Kaygorodov 2007). It also treats the
related $\tfrac{1}{2}$-derivations on Lie algebras that encode right
multiplication in transposed Poisson algebras
\cite[Theorems~2 and~3, pp.~356--358]{Filippov2000DeltaAlternative},
\cite[Theorems~2.5 and~4.4]{KacKaygorodov2007Delta},
\cite[Definition~5, p.~4]{FerreiraKaygorodovLopatkin21}.
Section~\ref{sec:nonassociative:poisson} gives the definition of Li of a
Poisson derivation and the theorem that every Poisson derivation of the
symplectic Poisson algebra $S_n$ is inner
\cite[Definitions, p.~658]{AlbertiArestiLi2018Poisson},
\cite[Theorem~2.1, pp.~658--659]{AlbertiArestiLi2018Poisson}.
Section~\ref{sec:nonassociative:hom-conformal} compares two results: the
Hom-Lie algebras of Hartwig, Larsson and Silvestrov, and the triviality
of such structures on simple Lie algebras proved by Jin and Li
\cite[Definition~4, p.~18]{HartwigLarssonSilvestrov2006},
\cite[Proposition~2.2, p.~1400]{JinLi2008HomLie}. It closes with the
definition of a conformal derivation of a Lie conformal algebra of Fan,
Hong and Su \cite[Definitions~2.1 and~2.7, pp.~2--3]{FanHongSu2019}.
Section~\ref{sec:nonassociative:n-ary} closes the chapter with the
derivations of the ternary Malcev algebra $M_8$ and with the
$(n+1)$-ary derivations of simple Mal'tsev algebras
\cite[Theorem~2.3, p.~3598]{PojidaevSarahva2006},
\cite[Theorem~5, p.~582]{Kaygorodov2014n1aryMaltsev}.
\hbfigure{15-nonassociative/constructions}{From associative to Lie and Jordan algebras, and from alternative to Malcev algebras. A derivation of $A$ is also a derivation of $A^{-}$ and of $A^{+}$.}{fig:nonassoc:constructions}
